% SPDX-License-Identifier: CC-BY-4.0
\documentclass[11pt,letterpaper]{article}

\usepackage[T1]{fontenc}
\usepackage{amsmath,amssymb}
\usepackage{tikz}
\usepackage{pgfplots}
\pgfplotsset{compat=1.18}
\usepackage{algorithm}
\usepackage{algpseudocode}
\usepackage{booktabs}
\usepackage{enumitem}
\usepackage[hidelinks]{hyperref}

\newcommand{\N}{\mathbb{N}}
\newcommand{\Z}{\mathbb{Z}}
\newcommand{\Q}{\mathbb{Q}}
\newcommand{\R}{\mathbb{R}}
\newcommand{\C}{\mathbb{C}}
\newcommand{\Map}[3]{#1 \colon #2 \to #3}
\newcommand{\Mod}[1]{#1\text{-}\mathrm{Mod}}
\newcommand{\GroupHom}{\mathrm{Hom}_{\mathbf{Grp}}}
\newcommand{\RingHom}{\mathrm{Hom}_{\mathbf{Ring}}}

\newtheorem{theorem}{Theorem}
\newtheorem{lemma}[theorem]{Lemma}
\newtheorem{corollary}[theorem]{Corollary}
\newtheorem{definition}[theorem]{Definition}
\newtheorem{remark}[theorem]{Remark}
\newenvironment{proof}{\noindent\textit{Proof.}}{\hfill\(\square\)}

\title{Exact product of a signed division}
\author{}
\date{\today}

\begin{document}
\maketitle

\begin{abstract}
If a nonzero integer \(d\) divides an integer \(a\), the truncating quotient
\(a/d\) is the divisibility witness, and therefore \(d\cdot(a/d)=a\). Absolute
values multiply: \(|d|\cdot|a/d|=|a|\), including when \(d\) is negative. Both
identities apply to the affine state at a first hit. Each one is an equality.
No nonnegative remainder is discarded, and the triangle inequality is not used.
\end{abstract}

\section{Division on exact multiples}

Throughout, \(d\) and \(a\) denote integers. The symbol \(d\mid a\) means that
there exists an integer \(k\) with \(a=dk\). For \(d\neq 0\), truncating
division writes every integer in the form \(a=dq+r\) with \(|r|<|d|\). On an
exact multiple the remainder is forced to vanish, so the quotient is the
witness itself.

\begin{lemma}
\label{lem:cancel}
Let \(d,k,k'\in\Z\). If \(d\neq 0\) and \(dk=dk'\), then \(k=k'\).
\end{lemma}

\begin{proof}
The hypothesis gives \(d(k-k')=0\). Suppose \(k-k'\neq 0\). Then
\(|k-k'|\ge 1\). Also \(|d|\ge 1\), because \(d\neq 0\). Hence
\[
|d(k-k')|=|d|\cdot|k-k'|\ge 1,
\]
so \(d(k-k')\neq 0\), a contradiction. Therefore \(k-k'=0\).
\end{proof}

\begin{lemma}
\label{lem:quotient-witness}
Let \(d,a,k\in\Z\) with \(d\neq 0\) and \(a=dk\). The truncating quotient
satisfies \(a/d=k\).
\end{lemma}

\begin{proof}
Truncating division supplies integers \(q\) and \(r\) with
\[
a=dq+r,\qquad |r|<|d|.
\]
The pair \(q=k\) and \(r=0\) is one such representation, since
\(a=dk+0\) and \(|0|<|d|\). For this same representation, or for any other
pair with \(|r|<|d|\),
\[
d(k-q)=r.
\]
Taking absolute values, \(|d|\cdot|k-q|=|r|<|d|\). If \(k\neq q\), then
\(|k-q|\ge 1\), so \(|d|\le |r|\), contradicting \(|r|<|d|\). Thus \(q=k\) and
\(r=0\). The truncating quotient is this \(q\), so \(a/d=k\).
\end{proof}

The same conclusion holds for a negative divisor. The comparison uses
\(|d|\), not the sign of \(d\), and Lemma~\ref{lem:cancel} does not require
\(d>0\).

\begin{theorem}
\label{thm:return-product}
Let \(d,a\in\Z\). If \(d\neq 0\) and \(d\mid a\), then
\[
d\cdot(a/d)=a.
\]
\end{theorem}

\begin{proof}
By the definition of divisibility, there exists \(k\in\Z\) with \(a=dk\).
Lemma~\ref{lem:quotient-witness} gives \(a/d=k\). Substituting the witness,
\[
d\cdot(a/d)=dk=a.
\]
\end{proof}

\begin{lemma}
\label{lem:abs-mul}
For all \(x,y\in\Z\), \(|xy|=|x|\cdot|y|\).
\end{lemma}

\begin{proof}
The identity is checked on the sign of each factor, including zero. If \(x=0\)
or \(y=0\), then both sides are zero. Now assume \(x\neq 0\) and \(y\neq 0\).

If \(x>0\) and \(y>0\), then \(|x|\cdot|y|=xy=|xy|\).

If \(x>0\) and \(y<0\), then \(|y|=-y\), so
\[
|x|\cdot|y|=x(-y)=-(xy)=|xy|,
\]
where the last step uses \(xy<0\).

If \(x<0\) and \(y>0\), then \(|x|=-x\), so
\[
|x|\cdot|y|=(-x)y=-(xy)=|xy|,
\]
where the last step uses \(xy<0\).

If \(x<0\) and \(y<0\), then \(xy>0\), \(|x|=-x\), and \(|y|=-y\), so
\[
|x|\cdot|y|=(-x)(-y)=xy=|xy|.
\]
In every case the two sides are equal. The argument does not replace \(|xy|\)
by \(|x|+|y|\) or by any larger quantity.
\end{proof}

\begin{theorem}
\label{thm:return-abs-product}
Let \(d,a\in\Z\). If \(d\neq 0\) and \(d\mid a\), then
\[
|d|\cdot|a/d|=|a|.
\]
\end{theorem}

\begin{proof}
Choose \(k\in\Z\) with \(a=dk\), as in Theorem~\ref{thm:return-product}.
Lemma~\ref{lem:quotient-witness} gives \(a/d=k\). Therefore
\[
|d|\cdot|a/d|=|d|\cdot|k|.
\]
Lemma~\ref{lem:abs-mul}, applied to the factors \(d\) and \(k\), yields
\(|d|\cdot|k|=|dk|\). Since \(a=dk\),
\[
|d|\cdot|a/d|=|a|.
\]
The four sign patterns of Lemma~\ref{lem:abs-mul} include every case in which
\(d\) is negative. No nonnegative term is removed from either side: the
remainder in Lemma~\ref{lem:quotient-witness} is \(0\), and the absolute-value
step is an equality.
\end{proof}

\begin{remark}
\label{rem:not-triangle}
Theorem~\ref{thm:return-abs-product} is not the estimate
\(|d\cdot(a/d)|\le |d|\cdot|a/d|\), nor the triangle inequality. For an exact
multiple those comparisons become equalities, because
\(d\cdot(a/d)=a\) and \(|dk|=|d|\cdot|k|\). A non-multiple is not assigned a
rounded quotient: if \(d\) does not divide \(a\), the hypothesis of both
theorems fails, and the identity \(d\cdot(a/d)=a\) is not claimed.
\end{remark}

\section{The affine state at a first hit}

Let \(q,c\in\Z\) and let \(j\in\N\), including zero. The offset sum is defined
by \(T_0(q)=0\) and \(T_{j+1}(q)=T_j(q)+q^j\).

\begin{definition}
\label{def:affine-state}
The affine state at time \(j\) is
\[
A_j(q,c,n)=q^j n+c\,T_j(q).
\]
\end{definition}

A parameter record consists of integers \(q\), \(c\), and \(d\) with
\(|d|\ge 2\). In particular \(d\neq 0\). The raw step sends an integer \(n\) to
\(n/d\) when \(d\mid n\), and to \(qn+c\) otherwise. The theorems below do not
iterate that step. They apply the product identities to one affine state.

\begin{definition}
\label{def:first-hit}
Let \(p\) be a parameter record with multiplier \(q\), offset \(c\), and
divisor \(d\). The time \(j\) is a first hit of the seed \(n\) when
\(d\mid A_j(q,c,n)\) and \(d\nmid A_i(q,c,n)\) for every natural number
\(i<j\).
\end{definition}

\begin{theorem}
\label{thm:return-product-hit}
Let \(p\) be a parameter record with divisor \(d\), and let \(n\in\Z\) and
\(j\in\N\). If \(j\) is a first hit of \(n\), and \(a=A_j(q,c,n)\) is the
affine state of that record at time \(j\), then
\[
d\cdot(a/d)=a.
\]
\end{theorem}

\begin{proof}
The first-hit hypothesis gives \(d\mid a\). The admissibility condition
\(|d|\ge 2\) gives \(d\neq 0\). Theorem~\ref{thm:return-product} applies to
this pair \((d,a)\) and yields \(d\cdot(a/d)=a\).

The clause that \(d\) divides no earlier affine state is part of the
first-hit hypothesis. The cancellation uses only \(d\mid a\) and \(d\neq 0\).
It does not induct on \(j\), and it does not rewrite any state at a time
other than \(j\).
\end{proof}

\begin{theorem}
\label{thm:return-abs-product-hit}
Under the hypotheses of Theorem~\ref{thm:return-product-hit},
\[
|d|\cdot|a/d|=|a|.
\]
\end{theorem}

\begin{proof}
As in Theorem~\ref{thm:return-product-hit}, \(d\mid a\) and \(d\neq 0\).
Theorem~\ref{thm:return-abs-product} gives the absolute-value equality. The
sign of \(d\) remains unrestricted beyond \(|d|\ge 2\), so a negative divisor
is included by Lemma~\ref{lem:abs-mul}.
\end{proof}

\begin{remark}
The two first-hit theorems recover the divisible affine state from its
quotient. They do not assert that a first hit exists, and they do not describe
the raw orbit at any time after \(j\). If \(d\) divides \(a\), the raw step at
that single state is the quotient \(a/d\), and Theorem~\ref{thm:return-product-hit}
says that multiplying by \(d\) returns \(a\). That is the whole claim.
\end{remark}

\end{document}
